% Generated by report_gen.py (hwde v1) - do not hand-edit.
\documentclass[11pt]{article}
\newcommand\hspaper{a4paper}
\usepackage[nodiagramkit]{housestyle}
\usepackage{longtable}
\usepackage{pdfpages}
\makeatletter
\@ifundefined{KV@Gin@artifact}{\define@key{Gin}{artifact}[]{}}{}
\makeatother
\setcounter{tocdepth}{1}
\hsslug{PCB-0021-A \textperiodcentered\ lipo-boost}
\renewcommand{\sectionmark}[1]{\markboth{#1}{}}
\hssection{\leftmark}
\begin{document}
\sloppy
\hstitleblock{PCB-0021-A \textperiodcentered\ hwde v1 pipeline \textperiodcentered\ generated 2026-10-01 05:51}%
{lipo-boost \textemdash{} Design Document}%
{Phase P2; no gates recorded yet.}
\tableofcontents

\section{Board and Run Metadata}
{\small
\begin{longtable}{lp{11cm}}
\hstoprule
\hshead{Field} & \hshead{Value} \\
\hline
\endhead
board & lipo-boost \\
part number & PCB-0021-A \\
workspace & \texttt{PCB-0021-A\_lipo-boost} \\
phase & P2 \\
gate status & no gates recorded yet \\
state created & 2026-10-01T05:34:53 \\
state updated & 2026-10-01T05:51:33 \\
\hline
\end{longtable}
}
\section{Overview}
\subsection*{lipo-boost brief (PCB-0021-A)}
I want a little board that runs a 5 V USB gadget off a single-cell LiPo and charges the cell from USB.

\begin{itemize}
\item USB in (micro-B or C) charges a 1S LiPo on a JST-PH 2-pin connector, about 1 A max.
\item Load sharing: when USB is plugged in, the 5 V output keeps running while the cell charges.
\item 5 V / 1 A out on a USB-A socket, from a synchronous boost.
\item An enable pin to turn the output off, plus a low-battery LED and charge / done LEDs.
\item Small, cheap, JLCPCB-assembled, 2 layers is fine.
\end{itemize}
Reference: this is a from-scratch recreation of Adafruit's PowerBoost 1000C (open hardware, CC BY-SA). Design it from this brief, then compare the result with Adafruit's published design.

\section{Requirements}
\subsection*{Requirements - lipo-boost (PCB-0021-A)}
\subsubsection*{1. Function}
A small 2-layer board that runs a 5 V / 1 A USB gadget from a single-cell LiPo and charges that cell from USB. USB input charges the cell through a linear charger, a power-path (load sharing) keeps the 5 V output running while USB is present, and a synchronous boost converter makes 5 V from the cell otherwise. Recreation of the open-hardware PowerBoost 1000C class of design, built from this brief. No mode token: normal scope, geometry is an output of design within the soft limits below.

\subsubsection*{2. Interfaces}
\begin{itemize}
\item USB input: USB micro-B receptacle (see Answers Q1), 5 V charge input only, no data.
\item Battery: JST-PH 2-pin, 1S LiPo, polarity per Adafruit/JST-PH convention (see Answers Q3).
\item Output: USB-A receptacle, 5 V / 1 A, power only (D+/D- left unconnected or shorted per Answers Q9).
\item EN input: 2-pin header or pad; pulling low turns the 5 V output off.
\item LEDs: charge (CHG), done (DONE), low-battery (LBO).
\end{itemize}
\subsubsection*{3. Power}
\begin{itemize}
\item Input: USB 5 V (4.5-5.5 V), up to \textasciitilde{}1.5 A draw (1 A charge plus load).
\item Cell: 1S LiPo 3.0-4.2 V, assumed protected (own protection circuit).
\item Charge current: 1 A default (Answers Q2), CC/CV 4.2 V, termination \textasciitilde{}10\% of set current.
\item Output: 5 V at 1 A continuous from a synchronous boost; at 3.2 V cell that is \textasciitilde{}2 A input, GUESS peak switch current \textasciitilde{}3 A.
\item Load sharing: output stays up when USB present (ideal-diode / power-path, Answers Q4).
\item Low-battery threshold \textasciitilde{}3.2 V (Answers Q5). Rail budget guesses: LED current \textasciitilde{}1-2 mA each.
\end{itemize}
\subsubsection*{4. Environment}
Indoor, 0-40 C, no enclosure specified, no ingress or vibration requirement. ASSUMED.

\subsubsection*{5. Size \& mounting}
\begin{itemize}
\item Outline: RELAXABLE (no HARD cap); target roughly 40 x 30 mm or smaller.
\item Mounting: none required; 2 bare M2.5 holes optional (ASSUMED none, Answers Q10).
\item Height: no limit; USB-A and JST dominate.
\end{itemize}
\subsubsection*{6. Quantity \& budget}
Quantity 5 boards (ASSUMED, prototype). Target unit cost under about USD 8 assembled, 2 layers.

\subsubsection*{7. Assembly}
JLCPCB PCBA, single-sided assembly preferred (top), prefer basic/preferred library parts; through-hole USB-A and JST may be hand-fitted or JLC THT.

\subsubsection*{8. Compliance/safety flags}
\begin{itemize}
\item Battery: 1S LiPo, 4.2 V charge, up to 1 A charge current: YES. Cell assumed protected; board gives no cell protection beyond charger limits.
\item High current: boost input \textasciitilde{}2-3 A peak: below the 3 A flag in steady state but near it; layout must use wide power traces.
\item NTC thermistor: absent (Answers Q6).
\item No mains, no motors, no \textgreater{}30 V, no RF.
\item Charging a LiPo unattended carries fire risk; board is a prototype, not certified.
\end{itemize}
\subsubsection*{9. Open questions}
All answered below as PROVISIONAL delegate answers (owner delegated all choices, unattended run). No question remains for the owner: 1. Micro-B or USB-C input? 2. Charge current? 3. Cell protected / polarity? 4. Power-path type? 5. Low-battery threshold? 6. NTC? 7. EN polarity/default? 8. Output USB-A data lines? 9. Output current limit? 10. Mounting holes?

\subsubsection*{Answers}
All PROVISIONAL (delegate answers, owner not asked). 1. Input connector: micro-B, cheaper and reference-like; no CC resistors needed. PROVISIONAL. 2. Charge current: 1 A default, set by a programming resistor on the charger. PROVISIONAL. 3. Cell: assumed to have its own protection circuit (PCM); JST-PH polarity per Adafruit convention (pin 1 +, to be stated in the design and silk-marked). PROVISIONAL. 4. Load sharing: ideal-diode/power-path via a P-MOSFET or Schottky from the charger/cell node, choice left to architecture. PROVISIONAL. 5. Low-battery LED lights at about 3.2 V cell. PROVISIONAL. 6. NTC: optional and absent (charger temperature pin tied to a fixed resistor/divider as in the reference). PROVISIONAL. 7. EN: active-high, pulled up on board so output is on by default; pull to GND turns it off. PROVISIONAL. 8. USB-A D+/D-: left floating, power only. PROVISIONAL. 9. Output current: 1 A rated; boost limit about 2 A; no extra eFuse. PROVISIONAL. 10. Mounting holes: none. PROVISIONAL.

\section{Architecture}
\subsection*{Decisions of record (state.json)}
{\small
\begin{longtable}{lp{5.2cm}p{7.2cm}}
\hstoprule
\hshead{Phase} & \hshead{What} & \hshead{Why} \\
\hline
\endhead
P2 & PROVISIONAL: 2L JLC2313\_1.6, single-sided SMT (USB-A THT); 2 oz fallback & small 5 V/1 A board; unattended delegate \\
P2 & PROVISIONAL: one schematic agent, two sheets charger + boost & small board \\
P2 & PROVISIONAL: charger BQ24074RGTR, boost TPS61023DRLR & OUT regulated 4.4 V so boost never passes USB through; TPS61023 cheap, 3.7 A limit, JLC stock \\
P2 & PROVISIONAL: USB-A D+/D- shorted by 0R (DCP) & BC1.2 DCP; one part; Apple dividers fail BC1.2 \\
P2 & PROVISIONAL: Vout 4.95 V (732k/100k) & stays in 4.75-5.25 V across tolerance \\
P2 & PROVISIONAL: ICHG 1 A (RISET), ILIM 1.5 A, TS fixed 10k, no NTC & brief 1 A; same as reference; cells \textless{}1000 mAh need RISET raised \\
P2 & PROVISIONAL: DONE LED via one NMOS gate=CHG source=PGOOD; LOW-BATT 3.2 V detector drives LED + LBO pin & BQ24074 has no done pin; TPS61023 has no LBO \\
P2 & PROVISIONAL: EN 100k pull-up (on by default); header GND,BAT,VS,5V,LBO,EN,GND & matches reference pinout \\
P2 & PROVISIONAL: no PTC/eFuse on 5V, no reverse-battery FET, silk polarity warning only & delegate answer Q9; same as reference; safety-relevant for owner review \\
P2 & research task block-boost-1 closed over 2 unverified draft record(s): boost-fb-divider-and-feedforward, boost-thermal-copper-and-vias & --accept-drafts: the slot ships without this knowledge \\
P2 & research task block-linear-charger-power-path-1 closed over 4 unverified draft record(s): linear-charger-ilim-iset-programming, linear-charger-pin-caps-ts-and-leds, linear-charger-single-point-ground-and-current-paths, linear-charger-thermal-dissipation-and-pad & --accept-drafts: the slot ships without this knowledge \\
P2 & Closed both P2 research tasks --accept-drafts (6/12 verified, 6 refuted drafts stay draft and do not inject) & budget; refutations are numeric corrections carried as the decision below, not wrong rules \\
P2 & Design corrections from the second read: TPS61023 VREF 580/595/610 mV (PWM) for the FB divider tolerance; BQ24074 size ILIM/ISET with KILIM max 1720 (I = K/R); bq24074 VIN range 4.35-10.2 V (VOVP 10.2-10.8 V); TPS61023 RthJA 142.7 C/W -\textgreater{} \textasciitilde{}0.3-0.55 W loss is a 43-78 C rise on 2L, so give U2 a GND pour and vias; ISET pads need no dedicated vias & second reader refutations, research/records \\
P2 & designing under coverage gap: block:B1 linear-charger-power-path and block:B2 boost (checklists + 6 records still draft after second read) & first research pass on these topologies; verified records (boost hot-loop, sw-node, inductor, caps, enable; charger DPM) guide layout by judgement \\
\hline
\end{longtable}
}
\subsection*{Sheet plan}
\subsection*{Sheet plan - PCB-0021-A (one schematic agent, two sheets; PROVISIONAL)}
\subsubsection*{Sheet \texttt{charger} (root file lipo-boost.kicad\_sch, local sheet)}
Blocks: U1 BQ24074, J1 micro-B, J2 JST-PH, D1/D2/D3 LEDs, Q1, U3 detector, TS/ISET/ILIM/TMR/ITERM resistors, USB-in/BAT/SYS caps. Refdes: U1,U3; Q1; D1-D3; J1,J2; R1-R19; C1-C9; \#PWR 100+. Nets: local /CHG, /PGOOD, /ISET, /ILIM, /TMR, /ITERM, /TS, /DONE\_K, /LB\_K; power VBUS, VBAT, +SYS, GND. Interface (hier pins) to \texttt{boost}: +SYS (out), VBAT (out), GND, /LBO (out), /EN (in from header).

\subsubsection*{Sheet \texttt{boost}}
Blocks: U2 TPS61023, L1, FB divider, EN pull-up, Cin/Cout, J4 USB-A (+R\_DCP 0 R), J3 header. Refdes: U2; L1; R20-R39; C10-C29; J3,J4; \#PWR 200+. Nets: local /SW, /FB, /EN, /USB\_DP, /USB\_DM (joined by R\_DCP); power +SYS, VBAT, +5V, GND. Interface to \texttt{charger}: +SYS (in), VBAT (in), GND, /LBO (in, to header), /EN. Placement groups at P6: cin\_u2\_l1 (C10, U2, L1 hot loop), cout (C11,C12), u1\_thermal (U1 + C\_out/C\_in), fb (R\_FB1,R\_FB2 away from L1).

Full architecture narrative (not inlined; see the artifact index): \texttt{architecture/blocks.md}, \texttt{architecture/power\_tree.md}, \texttt{architecture/stackup.md}.
\section{Schematic}
\emph{Pending \textemdash{} produced at P4.}
\section{Layout}
\emph{Pending \textemdash{} produced at P6.}
\section{Verification}
\emph{Pending \textemdash{} produced at P8.}
\section{DFM and Fabrication}
\emph{Pending \textemdash{} produced at P9.}
\section{Run Record (Appendix)}
\subsection*{Phase digests}
\paragraph*{P0}
\begin{quote}\small\ttfamily\raggedright
\# P0 digest\\
- requirements.md by requirements-analyst (sonnet); check\_requirements pass.\\
- 1S LiPo JST-PH, micro-B in, 1 A CC/CV charge to 4.2 V, power-path load sharing, sync boost 5 V/1 A to USB-A.\\
- EN active-high pulled up (on by default); charge/done/low-batt LEDs; low-batt \textasciitilde{}3.2 V.\\
- 2L JLC PCBA single-sided, outline relaxable, no mounting holes.\\
- 10 OPEN answered as PROVISIONAL delegate (unattended run): cell protected, NTC absent, no eFuse, D+/D- policy left to P2.\\
- Risks: boost peak input \textasciitilde{}2-3 A at low cell; power-path part; JST-PH polarity silk.
\end{quote}
\paragraph*{P1}
\begin{quote}\small\ttfamily\raggedright
\# P1 digest\\
- One combined sonnet researcher (power-architect + scout + reference-design) to save spend: research/power.md, parts-scout.md, reference-adafruit-powerboost1000c.md.\\
- Charger/power-path: BQ24074RGTR C54313 (OUT regulated 4.4 V, so boost never passes USB through); fallback BQ24072RGTR C140288.\\
- Boost: TPS61023DRLR C919459 (SOT-563, 3.7 A valley limit, \textasciitilde{}94 \%); inductor 1 uH Isat \textgreater{}= 4 A, to verify.\\
- Reference: Adafruit PowerBoost 1000C, CC BY-SA 3.0 Unported; MCP73871 + TPS61090.\\
- Gaps for P2/P3: low-battery detector (TPS61023 has no LBO), DONE LED (BQ24074 has CHG + PGOOD only), TS tied to 10k (no NTC, same as reference), charger heat \textasciitilde{}1.1 W.\\
- Prices/stock from LCSC web 2026-10-01 (parts\_search offline); re-verify at P3.
\end{quote}
\paragraph*{P2}
\emph{(no digest recorded)}
\subsection*{Gate attempt history}
\emph{(no entries)}
\subsection*{Issues}
\emph{(no entries)}
\subsection*{Human checkpoints}
{\small
\begin{longtable}{lllp{2.6cm}p{6.4cm}}
\hstoprule
\hshead{Id} & \hshead{Phase} & \hshead{Status} & \hshead{When} & \hshead{Note} \\
\hline
\endhead
1 & P2 & not reached &  &  \\
2 & P4 & not reached &  &  \\
3 & P6 & not reached &  &  \\
4 & P8 & not reached &  &  \\
5 & P10 & not reached &  &  \\
\hline
\end{longtable}
}
\subsection*{Generation cost}
\emph{Not recorded (no reports/cost.json).}
\section{Artifact Index}
{\small
\begin{longtable}{p{9cm}rp{4cm}}
\hstoprule
\hshead{File} & \hshead{Size} & \hshead{Note} \\
\hline
\endhead
\texttt{state.json} & 9,104 B &  \\
\texttt{brief/brief.md} & 705 B &  \\
\texttt{requirements.md} & 3,999 B &  \\
\texttt{architecture/blocks.md} & 2,593 B &  \\
\texttt{architecture/constraints.json} & 2,411 B &  \\
\texttt{architecture/power\_tree.md} & 728 B &  \\
\texttt{architecture/sheets.md} & 996 B &  \\
\texttt{architecture/stackup.md} & 881 B &  \\
\hline
\end{longtable}
}
\end{document}
